\documentclass[11pt]{article}
\usepackage[utf8]{inputenc}
\usepackage[T1]{fontenc}
\usepackage{booktabs}
\usepackage{amsmath}
\usepackage{geometry}
\geometry{margin=2.5cm}

\title{Causal Fairness Analysis Report}
\author{Generated by FairMind and an LLM}
\date{2026-08-22}

\begin{document}
\maketitle

\section{Analysis setup}

\begin{tabular}{ll}
\toprule
Dataset & student\_por \\
Protected attribute ($X$) & S2\_sex \\
Comparison & F $\to$ M \\
Outcome ($Y$) & T\_grade (1) \\
Mediator ($W$) & studytime \\
Confounder ($Z$) & Medu \\
\bottomrule
\end{tabular}

\section{Causal effects (FairMind, exact values)}

The five effects below are computed by FairMind through exact inference on the
fitted Bayesian network. These values are final and must not be recomputed or
altered.

\begin{tabular}{lr}
\toprule
Effect & Value \\
\midrule
Total Variation (TV) & -0.059600 \\
Total Effect (TE) & -0.079400 \\
Spurious Effect (SE) & 0.019800 \\
Direct Effect (DE) & -0.034800 \\
Indirect Effect (IE, decomposition form: TE = DE + IE) & -0.044600 \\
\bottomrule
\end{tabular}

\section{Qualitative interpretation}

\subsection{Total effect and spurious component (TV, TE, SE)}

The total disparity between the two groups is -0.0596, indicating a negative impact on the protected group. After removing the confounding effect of the mediator, the genuine causal effect (TE) is -0.0794, which is larger in magnitude than the total variation, suggesting that a significant portion of the observed disparity is due to genuine causal factors. The spurious effect (SE) of 0.0198 is driven by the confounder rather than the protected attribute itself, contributing to the observed difference but not as a direct causal effect.

\subsection{Direct effect (DE)}

The direct effect (DE) is -0.0348, indicating a negative impact on the protected group that is not mediated by the mediator variable. This suggests a direct discrimination against the protected group, though the magnitude is relatively small compared to the total effect.

\subsection{Indirect effect (IE)}

The indirect effect (IE) is -0.0446, which is of the same sign as the direct effect (DE). This means that the mediator channel amplifies the direct effect, as the magnitude of the total effect (TE) is larger than the magnitude of the direct effect (DE). Therefore, the indirect effect through the mediator does not mitigate the total effect but rather exacerbates it.

\section{Recap Questions}

\begin{enumerate}
\item Is there direct discrimination against the protected group? \textbf{Answer:} YES
\item Does the indirect effect through the mediator mitigate the total effect? \textbf{Answer:} NO
\item Does the observed total effect exceed the practical relevance threshold? \textbf{Answer:} YES
\item Does the spurious component (SE) contribute substantially to the observed difference? \textbf{Answer:} NO
\item Is the direct effect the dominant component compared to the indirect effect? \textbf{Answer:} NO
\end{enumerate}

\vspace{1em}
\noindent\footnotesize
The recap answers follow deterministic threshold rules applied to the values above: direct discrimination when $|\mathrm{DE}| > 0.05$; a mediator channel counted as negligible when $|\mathrm{IE}| < 0.005$; practical relevance when $|\mathrm{TV}| > 0.05$; a substantial spurious component when $|\mathrm{SE}| / |\mathrm{TV}| > 0.1$.

\end{document}